\section*{Problem 1}

\partlabel{(a)}
\textbf{(i) Isothermal:} $T=\text{constant}$.
% L3 notes, PDF p.1 (printed p.25), Eqs. (72)--(74): ideal gas and CPG.
% L1 notes, PDF p.13 (printed p.16), Eqs. (21)--(24): first law and work sign.
\begin{align*}
T_{2,\mathrm{iso}}&=T_1=\boxed{300~\mathrm{K}},
& pV&=nR_uT_1=\text{constant},
& \frac{V_{2,\mathrm{iso}}}{V_1}&=\frac{p_1}{p_2}=\frac18.
\end{align*}
For a calorically perfect gas, $\Delta E=n\bar c_v(T_2-T_1)=0$.
Thus the first law gives $0=Q-W_{\mathrm{by}}$.
The reversible boundary work is
% L1 notes PDF p.10 Eq. (9), p.13 Eq. (23); L3 PDF p.14 Eq. (164).
\begin{align*}
W_{\mathrm{by,iso}}
&=\int_{V_1}^{V_{2,\mathrm{iso}}}p\,dV
 =nR_uT_1\int_{V_1}^{V_{2,\mathrm{iso}}}\frac{dV}{V}\\
&=nR_uT_1\ln\frac{V_{2,\mathrm{iso}}}{V_1}
 =-(1)(8.31446)(300)\ln8
 =\boxed{-5.187~\mathrm{kJ}}.
\end{align*}

\smallskip
\noindent\textbf{(ii) Reversible adiabatic:} $Q=0$ and $S_{\mathrm{gen}}=0$,
so $\Delta S=0$.
% L2 notes PDF p.2 Eq. (25): entropy balance. L3 PDF p.1 Eq. (75),
% p.14 Eqs. (167)--(170): reversible adiabatic / isentropic CPG relations.
\begin{align*}
pV^\gamma&=\text{constant},
& \frac{T_{2,\mathrm{ad}}}{T_1}
 &=\left(\frac{p_2}{p_1}\right)^{(\gamma-1)/\gamma},\\
T_{2,\mathrm{ad}}&=300(8)^{0.4/1.4}
 =\boxed{543.4~\mathrm{K}}.
\end{align*}

% L2 notes PDF p.6 (printed p.21), Eq. (40), converted to a molar basis.
% M denotes molar mass here; cp_bar = M cp and cv_bar = M cv by definition.
\begin{align*}
\bar c_p&=\bar c_v+R_u,
& \gamma&=\frac{\bar c_p}{\bar c_v},\\
(\gamma-1)\bar c_v&=R_u,
& \bar c_v&=\frac{R_u}{\gamma-1}
 =20.7862~\mathrm{J/(mol\,K)}.
\end{align*}
Apply the first law with $Q=0$:
% L1 PDF p.13 Eqs. (21), (24) + L3 PDF p.1 Eq. (74).
% The work relation below is derived by combining these sourced equations.
\begin{align*}
\Delta E&=Q-W_{\mathrm{by,ad}}=-W_{\mathrm{by,ad}},\\
W_{\mathrm{by,ad}}&=-\Delta E=-m\Delta e
 =-n\bar c_v(T_{2,\mathrm{ad}}-T_1)
 =\boxed{-5.060~\mathrm{kJ}}.
\end{align*}

\partlabel{(b)}
% L2 notes PDF p.9 (printed p.24), Eq. (68); multiply by n M to obtain total S.
% L1 notation table PDF p.2: specific vs molar quantities.
Use the ideal-gas entropy relation and convert totals to mass-specific values:
\begin{gather*}
\Delta S_{\mathrm{sys}}
 =n\left(\bar c_p\ln\frac{T_2}{T_1}-R_u\ln\frac{p_2}{p_1}\right),
\qquad \bar c_p=\frac{\gamma R_u}{\gamma-1},\\
m=nM_{\mathrm{O_2}}=(1~\mathrm{mol})(0.03200~\mathrm{kg/mol})=0.03200~\mathrm{kg},
\qquad \Delta e=\frac{\Delta E}{m},\quad q=\frac{Q}{m},\quad
\Delta s_{\mathrm{sys}}=\frac{\Delta S_{\mathrm{sys}}}{m}.
\end{gather*}


\smallskip
\noindent\textbf{(i) Isothermal:}
\begin{align*}
\Delta E&=0,\qquad Q=W_{\mathrm{by,iso}}=-5.187~\mathrm{kJ},\\
\Delta S_{\mathrm{sys}}&=n\left(\bar c_p\ln1-R_u\ln8\right)
 =-nR_u\ln8=-17.29~\mathrm{J/K},
\end{align*}
\[
\boxed{\Delta e=0},\qquad
\boxed{q=-162.1~\mathrm{kJ/kg}},\qquad
\boxed{\Delta s_{\mathrm{sys}}=-0.5403~\mathrm{kJ/(kg\,K)}}.
\]

\smallskip
\noindent\textbf{(ii) Reversible adiabatic:}
\begin{align*}
\Delta E&=n\bar c_v(T_{2,\mathrm{ad}}-T_1)=+5.060~\mathrm{kJ},
\qquad Q=0,\\
\Delta S_{\mathrm{sys}}
&=n\left[\frac{\gamma R_u}{\gamma-1}
 \frac{\gamma-1}{\gamma}\ln8-R_u\ln8\right]=0,
\end{align*}
\[
\boxed{\Delta e=+158.1~\mathrm{kJ/kg}},\qquad
\boxed{q=0},\qquad \boxed{\Delta s_{\mathrm{sys}}=0}.
\]

\partlabel{(c)}
\textbf{(i) Isothermal:} 
% L2 notes PDF p.2 (printed p.17), Eq. (25), integrated for a reversible reservoir.
% Heat received by the surroundings is minus the heat received by the gas.
\begin{align*}
Q_{\mathrm{surr}}&=-Q=+5.187~\mathrm{kJ},\\
\Delta S_{\mathrm{surr}}&=\frac{Q_{\mathrm{surr}}}{300~\mathrm{K}}
 =+17.29~\mathrm{J/K},\\
\Delta S_{\mathrm{universe}}&=\Delta S_{\mathrm{sys}}+\Delta S_{\mathrm{surr}}
 =-17.29+17.29=\boxed{0}.
\end{align*}
\textbf{(ii) Reversible adiabatic:} 
\[
\boxed{\Delta S_{\mathrm{sys}}=\Delta S_{\mathrm{surr}}
=\Delta S_{\mathrm{universe}}=0.}
\]
The ideal work source transfers energy without entropy. Both processes have
zero entropy generation because they are reversible.

\partlabel{(d)}
% Work input is defined as minus work by; volumes follow from L3 PDF p.1 Eq. (72).
\begin{gather*}
W_{\mathrm{in}}=-W_{\mathrm{by}},\qquad
\boxed{W_{\mathrm{in,iso}}=5.187~\mathrm{kJ}},\qquad
\boxed{W_{\mathrm{in,ad}}=5.060~\mathrm{kJ}},\\
V_{2,\mathrm{iso}}=\frac{nR_uT_1}{p_2}=3.118~\mathrm{L},\qquad
V_{2,\mathrm{ad}}=\frac{nR_uT_{2,\mathrm{ad}}}{p_2}=5.648~\mathrm{L}.
\end{gather*}
For these data, the adiabatic path requires $0.1268~\mathrm{kJ}$ less work
input. Its temperature rises, so it reaches $8~\mathrm{bar}$ at a larger
volume and stops sooner. The final pressures are the same, but the final states are different. Comparing the paths at the same final volume would give a different work comparison.
